%%%%%%%%%%%%%%%%%%%%%%%%%%%%%%%%%%%%%%%%%%%%%%%%%%%%%%%%%%%%%%%%%%%%%%%%%%%%%%%%
% CHAPTER 12
%%%%%%%%%%%%%%%%%%%%%%%%%%%%%%%%%%%%%%%%%%%%%%%%%%%%%%%%%%%%%%%%%%%%%%%%%%%%%%%%

\chapter{Frequency Response, Poles \& Zeros, Phase Delay and Group Delay}

\begin{tcolorbox}[colback=blue!5!white,colframe=blue!70!black,title=Learning Objectives]

After studying this chapter, the reader will be able to

\begin{itemize}
\item Define the frequency response of an LTI system.
\item Relate the frequency response to the z-transform.
\item Understand poles and zeros in the \(z\)-plane.
\item Sketch pole-zero plots.
\item Determine system behavior from pole locations.
\item Compute magnitude and phase responses.
\item Understand resonance and bandwidth.
\end{itemize}

\end{tcolorbox}

%%%%%%%%%%%%%%%%%%%%%%%%%%%%%%%%%%%%%%%%%%%%%%%%%%%%%%%%%%%%%%%%%%%%%%%%%%%%%%%%

\section{Introduction}

An LTI system is completely characterized by its impulse response \(h[n]\).
However, in signal processing and communication systems, we are often interested
in how the system responds to different frequencies.

The \textbf{frequency response} tells us

\begin{itemize}
\item which frequencies are passed,
\item which frequencies are attenuated,
\item and how much phase shift each frequency experiences.
\end{itemize}

%%%%%%%%%%%%%%%%%%%%%%%%%%%%%%%%%%%%%%%%%%%%%%%%%%%%%%%%%%%%%%%%%%%%%%%%%%%%%%%%

\section{System Function}

For a discrete-time LTI system,

\[
H(z)=\frac{Y(z)}{X(z)}.
\]

If

\[
H(z)=
\frac{b_0+b_1z^{-1}+\cdots+b_Mz^{-M}}
{a_0+a_1z^{-1}+\cdots+a_Nz^{-N}},
\]

then

\begin{itemize}
\item the numerator determines the \textbf{zeros},
\item the denominator determines the \textbf{poles}.
\end{itemize}

%%%%%%%%%%%%%%%%%%%%%%%%%%%%%%%%%%%%%%%%%%%%%%%%%%%%%%%%%%%%%%%%%%%%%%%%%%%%%%%%

\section{Frequency Response}

The frequency response is obtained by evaluating \(H(z)\) on the unit circle:

\[
\boxed{
H(e^{j\omega})=
H(z)\Big|_{z=e^{j\omega}}
}
\]

This is valid only if the unit circle lies inside the ROC.

%%%%%%%%%%%%%%%%%%%%%%%%%%%%%%%%%%%%%%%%%%%%%%%%%%%%%%%%%%%%%%%%%%%%%%%%%%%%%%%%

\section{Magnitude and Phase}

Any complex frequency response can be written as

\[
\boxed{
H(e^{j\omega})=
|H(e^{j\omega})|e^{j\angle H(e^{j\omega})}
}
\]

where

\[
|H(e^{j\omega})|
\]

is the magnitude response and

\[
\angle H(e^{j\omega})
\]

is the phase response.

%%%%%%%%%%%%%%%%%%%%%%%%%%%%%%%%%%%%%%%%%%%%%%%%%%%%%%%%%%%%%%%%%%%%%%%%%%%%%%%%

\section{Example 1}

Consider

\[
H(z)=
\frac{1}{1-0.5z^{-1}}.
\]

%%%%%%%%%%%%%%%%%%%%%%%%%%%%%%%%%%%%%%%%%%%%%%%%%%%%%%%%%%%%%%%%%%%%%%%%%%%%%%%%

\subsection{Frequency Response}

Substitute \(z=e^{j\omega}\):

\[
H(e^{j\omega})=
\frac{1}{1-0.5e^{-j\omega}}.
\]

%%%%%%%%%%%%%%%%%%%%%%%%%%%%%%%%%%%%%%%%%%%%%%%%%%%%%%%%%%%%%%%%%%%%%%%%%%%%%%%%

\subsection{Expand}

\[
1-0.5e^{-j\omega}
=
1-0.5\cos\omega+j0.5\sin\omega.
\]

%%%%%%%%%%%%%%%%%%%%%%%%%%%%%%%%%%%%%%%%%%%%%%%%%%%%%%%%%%%%%%%%%%%%%%%%%%%%%%%%

\subsection{Magnitude}

\[
\boxed{
|H(e^{j\omega})|=
\frac{1}
{\sqrt{1.25-\cos\omega}}
}
\]

%%%%%%%%%%%%%%%%%%%%%%%%%%%%%%%%%%%%%%%%%%%%%%%%%%%%%%%%%%%%%%%%%%%%%%%%%%%%%%%%

\subsection{Phase}

\[
\boxed{
\angle H(e^{j\omega})=
-\tan^{-1}
\left(
\frac{0.5\sin\omega}{1-0.5\cos\omega}
\right)
}
\]

%%%%%%%%%%%%%%%%%%%%%%%%%%%%%%%%%%%%%%%%%%%%%%%%%%%%%%%%%%%%%%%%%%%%%%%%%%%%%%%%

\section{Poles and Zeros}

%%%%%%%%%%%%%%%%%%%%%%%%%%%%%%%%%%%%%%%%%%%%%%%%%%%%%%%%%%%%%%%%%%%%%%%%%%%%%%%%

\subsection{Zero}

A value of \(z\) that makes

\[
H(z)=0.
\]

%%%%%%%%%%%%%%%%%%%%%%%%%%%%%%%%%%%%%%%%%%%%%%%%%%%%%%%%%%%%%%%%%%%%%%%%%%%%%%%%

\subsection{Pole}

A value of \(z\) that makes

\[
H(z)\to\infty.
\]

%%%%%%%%%%%%%%%%%%%%%%%%%%%%%%%%%%%%%%%%%%%%%%%%%%%%%%%%%%%%%%%%%%%%%%%%%%%%%%%%

\subsection{Example}

\[
H(z)=
\frac{1-z^{-1}}
{1-0.5z^{-1}}.
\]

Zero:

\[
z=1.
\]

Pole:

\[
z=0.5.
\]

%%%%%%%%%%%%%%%%%%%%%%%%%%%%%%%%%%%%%%%%%%%%%%%%%%%%%%%%%%%%%%%%%%%%%%%%%%%%%%%%

\section{Pole-Zero Plot}

\begin{center}
\begin{tikzpicture}[scale=2]
\draw[->] (-1.5,0) -- (1.5,0) node[right] {Re};
\draw[->] (0,-1.2) -- (0,1.2) node[above] {Im};

\draw (0,0) circle (1);

\draw[thick] (1,0) node[circle,draw,inner sep=1.5pt] {};
\draw[thick] (0.5,0) node[cross out,draw,inner sep=2pt] {};

\node[below] at (1,0) {$z=1$};
\node[below] at (0.5,0) {$z=0.5$};
\end{tikzpicture}
\end{center}

Convention:

\begin{itemize}
\item \(\circ\): zero,
\item \(\times\): pole.
\end{itemize}

%%%%%%%%%%%%%%%%%%%%%%%%%%%%%%%%%%%%%%%%%%%%%%%%%%%%%%%%%%%%%%%%%%%%%%%%%%%%%%%%

\section{Effect of Zeros}

A zero on the unit circle creates a frequency null.

If

\[
z=e^{j\omega_0},
\]

then

\[
H(e^{j\omega_0})=0.
\]

The system completely rejects frequency \(\omega_0\).

%%%%%%%%%%%%%%%%%%%%%%%%%%%%%%%%%%%%%%%%%%%%%%%%%%%%%%%%%%%%%%%%%%%%%%%%%%%%%%%%

\section{Example: Notch at \(\omega=0\)}

For

\[
H(z)=1-z^{-1},
\]

the zero is at

\[
z=1=e^{j0}.
\]

Hence,

\[
H(e^{j0})=0.
\]

This removes the DC component.

%%%%%%%%%%%%%%%%%%%%%%%%%%%%%%%%%%%%%%%%%%%%%%%%%%%%%%%%%%%%%%%%%%%%%%%%%%%%%%%%

\section{Effect of Poles}

Poles close to the unit circle create peaks in the magnitude response.

%%%%%%%%%%%%%%%%%%%%%%%%%%%%%%%%%%%%%%%%%%%%%%%%%%%%%%%%%%%%%%%%%%%%%%%%%%%%%%%%

\subsection{Resonance}

If a pole is at

\[
z=re^{j\omega_0},
\qquad r\lesssim1,
\]

then the system exhibits resonance near \(\omega_0\).

%%%%%%%%%%%%%%%%%%%%%%%%%%%%%%%%%%%%%%%%%%%%%%%%%%%%%%%%%%%%%%%%%%%%%%%%%%%%%%%%

\section{Example: Resonant Pole}

\[
H(z)=
\frac{1}{1-0.95e^{j\pi/4}z^{-1}}.
\]

The pole is very close to the unit circle, so the magnitude response has a sharp
peak near

\[
\omega=\frac{\pi}{4}.
\]

%%%%%%%%%%%%%%%%%%%%%%%%%%%%%%%%%%%%%%%%%%%%%%%%%%%%%%%%%%%%%%%%%%%%%%%%%%%%%%%%

\section{Stability}

A causal system is stable if all poles lie strictly inside the unit circle:

\[
\boxed{
|p_i|<1
}
\]

%%%%%%%%%%%%%%%%%%%%%%%%%%%%%%%%%%%%%%%%%%%%%%%%%%%%%%%%%%%%%%%%%%%%%%%%%%%%%%%%

\subsection{Stable Pole}

\[
p=0.8
\]

%%%%%%%%%%%%%%%%%%%%%%%%%%%%%%%%%%%%%%%%%%%%%%%%%%%%%%%%%%%%%%%%%%%%%%%%%%%%%%%%

\subsection{Marginally Stable}

\[
p=1
\]

%%%%%%%%%%%%%%%%%%%%%%%%%%%%%%%%%%%%%%%%%%%%%%%%%%%%%%%%%%%%%%%%%%%%%%%%%%%%%%%%

\subsection{Unstable}

\[
p=1.2
\]

%%%%%%%%%%%%%%%%%%%%%%%%%%%%%%%%%%%%%%%%%%%%%%%%%%%%%%%%%%%%%%%%%%%%%%%%%%%%%%%%

\section{Frequency Response from Pole-Zero Distances}

For a point on the unit circle \(z=e^{j\omega}\),

\[
\boxed{
|H(e^{j\omega})|=
K
\frac{\prod_i |e^{j\omega}-z_i|}
{\prod_j |e^{j\omega}-p_j|}
}
\]

Interpretation:

\begin{itemize}
\item close to a zero \(\Rightarrow\) magnitude becomes small,
\item close to a pole \(\Rightarrow\) magnitude becomes large.
\end{itemize}

%%%%%%%%%%%%%%%%%%%%%%%%%%%%%%%%%%%%%%%%%%%%%%%%%%%%%%%%%%%%%%%%%%%%%%%%%%%%%%%%

\section{Low-Pass Filter Example}

\[
H(z)=
\frac{1}{1-0.5z^{-1}}.
\]

%%%%%%%%%%%%%%%%%%%%%%%%%%%%%%%%%%%%%%%%%%%%%%%%%%%%%%%%%%%%%%%%%%%%%%%%%%%%%%%%

\subsection{At \(\omega=0\)}

\[
|H(e^{j0})|=
\frac{1}{1-0.5}=2.
\]

%%%%%%%%%%%%%%%%%%%%%%%%%%%%%%%%%%%%%%%%%%%%%%%%%%%%%%%%%%%%%%%%%%%%%%%%%%%%%%%%

\subsection{At \(\omega=\pi\)}

\[
|H(e^{j\pi})|=
\frac{1}{1+0.5}=\frac23.
\]

Since low frequencies are amplified more than high frequencies,

\[
\boxed{
\text{Low-pass filter}
}
\]

%%%%%%%%%%%%%%%%%%%%%%%%%%%%%%%%%%%%%%%%%%%%%%%%%%%%%%%%%%%%%%%%%%%%%%%%%%%%%%%%

\section{High-Pass Filter Example}

\[
H(z)=1-z^{-1}.
\]

%%%%%%%%%%%%%%%%%%%%%%%%%%%%%%%%%%%%%%%%%%%%%%%%%%%%%%%%%%%%%%%%%%%%%%%%%%%%%%%%

\subsection{Magnitude}

\[
|H(e^{j\omega})|=
2\left|\sin\frac{\omega}{2}\right|.
\]

%%%%%%%%%%%%%%%%%%%%%%%%%%%%%%%%%%%%%%%%%%%%%%%%%%%%%%%%%%%%%%%%%%%%%%%%%%%%%%%%

\subsection{At \(\omega=0\)}

\[
|H|=0.
\]

%%%%%%%%%%%%%%%%%%%%%%%%%%%%%%%%%%%%%%%%%%%%%%%%%%%%%%%%%%%%%%%%%%%%%%%%%%%%%%%%

\subsection{At \(\omega=\pi\)}

\[
|H|=2.
\]

Therefore,

\[
\boxed{
\text{High-pass filter}
}
\]

%%%%%%%%%%%%%%%%%%%%%%%%%%%%%%%%%%%%%%%%%%%%%%%%%%%%%%%%%%%%%%%%%%%%%%%%%%%%%%%%

\section{Bandwidth}

The bandwidth is the frequency range over which the magnitude remains above the
\(-3\) dB level.

%%%%%%%%%%%%%%%%%%%%%%%%%%%%%%%%%%%%%%%%%%%%%%%%%%%%%%%%%%%%%%%%%%%%%%%%%%%%%%%%

\subsection{\(-3\) dB Condition}

\[
\boxed{
|H(e^{j\omega_c})|=
\frac{|H|_{\max}}{\sqrt{2}}
}
\]

The corresponding frequency \(\omega_c\) is the cutoff frequency.

%%%%%%%%%%%%%%%%%%%%%%%%%%%%%%%%%%%%%%%%%%%%%%%%%%%%%%%%%%%%%%%%%%%%%%%%%%%%%%%%

\section{First-Order Filter Cutoff}

For

\[
H(z)=
\frac{1}{1-az^{-1}},
\qquad 0<a<1,
\]

the cutoff frequency satisfies

\[
\boxed{
\cos\omega_c=
\frac{4a-a^2-1}{2a}
}
\]

As \(a\to1\),

\begin{itemize}
\item the bandwidth decreases,
\item selectivity increases.
\end{itemize}

%%%%%%%%%%%%%%%%%%%%%%%%%%%%%%%%%%%%%%%%%%%%%%%%%%%%%%%%%%%%%%%%%%%%%%%%%%%%%%%%

\section{Resonance and Quality Factor}

For a second-order system,

\[
H(z)=
\frac{1}
{1-2r\cos\omega_0 z^{-1}+r^2 z^{-2}},
\]

the parameter \(r\) controls the sharpness of resonance.

%%%%%%%%%%%%%%%%%%%%%%%%%%%%%%%%%%%%%%%%%%%%%%%%%%%%%%%%%%%%%%%%%%%%%%%%%%%%%%%%

\subsection{Approximate Relation}

\[
\boxed{
Q\approx\frac{1}{2(1-r)}
}
\]

Larger \(r\) gives higher \(Q\) and a narrower bandwidth.

%%%%%%%%%%%%%%%%%%%%%%%%%%%%%%%%%%%%%%%%%%%%%%%%%%%%%%%%%%%%%%%%%%%%%%%%%%%%%%%%

\section{Solved Example 12.1}

Find the poles and zeros of

\[
H(z)=
\frac{1-0.5z^{-1}}
{1-0.8z^{-1}}.
\]

%%%%%%%%%%%%%%%%%%%%%%%%%%%%%%%%%%%%%%%%%%%%%%%%%%%%%%%%%%%%%%%%%%%%%%%%%%%%%%%%

\subsection{Solution}

Zero:

\[
1-0.5z^{-1}=0
\Rightarrow z=0.5.
\]

Pole:

\[
1-0.8z^{-1}=0
\Rightarrow z=0.8.
\]

\[
\boxed{
\text{Zero at }0.5,\quad \text{Pole at }0.8
}
\]

%%%%%%%%%%%%%%%%%%%%%%%%%%%%%%%%%%%%%%%%%%%%%%%%%%%%%%%%%%%%%%%%%%%%%%%%%%%%%%%%

\section{Solved Example 12.2}

Determine whether the system is stable.

\[
H(z)=
\frac{1}{1-1.1z^{-1}}.
\]

%%%%%%%%%%%%%%%%%%%%%%%%%%%%%%%%%%%%%%%%%%%%%%%%%%%%%%%%%%%%%%%%%%%%%%%%%%%%%%%%

\subsection{Solution}

Pole:

\[
z=1.1.
\]

Since

\[
|1.1|>1,
\]

\[
\boxed{
\text{The system is unstable.}
}
\]

%%%%%%%%%%%%%%%%%%%%%%%%%%%%%%%%%%%%%%%%%%%%%%%%%%%%%%%%%%%%%%%%%%%%%%%%%%%%%%%%

\section{Solved Example 12.3}

Find the magnitude response of

\[
H(z)=1-z^{-1}.
\]

%%%%%%%%%%%%%%%%%%%%%%%%%%%%%%%%%%%%%%%%%%%%%%%%%%%%%%%%%%%%%%%%%%%%%%%%%%%%%%%%

\subsection{Solution}

\[
H(e^{j\omega})=1-e^{-j\omega}
\]

\[
=
e^{-j\omega/2}
\left(
e^{j\omega/2}-e^{-j\omega/2}
\right)
\]

\[
=
2je^{-j\omega/2}\sin\frac{\omega}{2}.
\]

Hence,

\[
\boxed{
|H(e^{j\omega})|=
2\left|\sin\frac{\omega}{2}\right|
}
\]

%%%%%%%%%%%%%%%%%%%%%%%%%%%%%%%%%%%%%%%%%%%%%%%%%%%%%%%%%%%%%%%%%%%%%%%%%%%%%%%%

\section{Solved Example 12.4}

Identify the filter type:

\[
H(z)=1+z^{-1}.
\]

%%%%%%%%%%%%%%%%%%%%%%%%%%%%%%%%%%%%%%%%%%%%%%%%%%%%%%%%%%%%%%%%%%%%%%%%%%%%%%%%

\subsection{Solution}

\[
H(e^{j\omega})=
1+e^{-j\omega}
=
2e^{-j\omega/2}\cos\frac{\omega}{2}.
\]

\[
|H(e^{j\omega})|=
2\left|\cos\frac{\omega}{2}\right|.
\]

At \(\omega=0\),

\[
|H|=2.
\]

At \(\omega=\pi\),

\[
|H|=0.
\]

\[
\boxed{
\text{Low-pass filter}
}
\]

%%%%%%%%%%%%%%%%%%%%%%%%%%%%%%%%%%%%%%%%%%%%%%%%%%%%%%%%%%%%%%%%%%%%%%%%%%%%%%%%

\section{Quick Check}

\begin{enumerate}
\item Frequency response is obtained by \(\boxed{z=e^{j\omega}}\).
\item Zero on unit circle \(\Rightarrow\) \(\boxed{\text{frequency null}}\).
\item Pole near unit circle \(\Rightarrow\) \(\boxed{\text{resonance}}\).
\item Stable causal system \(\Rightarrow\) \(\boxed{|p_i|<1}\).
\item \(1-z^{-1}\) is \(\boxed{\text{high-pass}}\).
\item \(1+z^{-1}\) is \(\boxed{\text{low-pass}}\).
\end{enumerate}

%%%%%%%%%%%%%%%%%%%%%%%%%%%%%%%%%%%%%%%%%%%%%%%%%%%%%%%%%%%%%%%%%%%%%%%%%%%%%%%%

\begin{tcolorbox}[title=One-Minute Revision,colback=blue!5]

\[
H(e^{j\omega})=H(z)|_{z=e^{j\omega}}
\]

\[
H(e^{j\omega})=
|H(e^{j\omega})|e^{j\angle H(e^{j\omega})}
\]

Zero: \(H(z)=0\).

Pole: denominator \(=0\).

Stable:

\[
|p_i|<1
\]

Low-pass:

\[
1+z^{-1}
\]

High-pass:

\[
1-z^{-1}
\]

Notch at \(\omega_0\):

\[
z=e^{\pm j\omega_0}
\]

\[
|H(e^{j\omega})|=
K
\frac{\prod |e^{j\omega}-z_i|}
{\prod |e^{j\omega}-p_i|}
\]

\end{tcolorbox}
%%%%%%%%%%%%%%%%%%%%%%%%%%%%%%%%%%%%%%%%%%%%%%%%%%%%%%%%%%%%%%%%%%%%%%%%%%%%%%%%
%%%%%%%%%%%%%%%%%%%%%%%%%%%%%%%%%%%%%%%%%%%%%%%%%%%%%%%%%%%%%%%%%%%%%%%%%%%%%%%%
\section{Phase Response}

The phase response of a system is

\[
\boxed{
\phi(\omega)=\angle H(e^{j\omega})
}
\]

It specifies the phase shift experienced by a sinusoidal component of frequency
\(\omega\).

%%%%%%%%%%%%%%%%%%%%%%%%%%%%%%%%%%%%%%%%%%%%%%%%%%%%%%%%%%%%%%%%%%%%%%%%%%%%%%%%

\section{Example}

Consider

\[
H(z)=1+z^{-1}.
\]

Substituting \(z=e^{j\omega}\),

\[
H(e^{j\omega})=
1+e^{-j\omega}.
\]

Factorize:

\[
H(e^{j\omega})=
2e^{-j\omega/2}\cos\frac{\omega}{2}.
\]

Hence,

\[
\boxed{
\phi(\omega)=-\frac{\omega}{2}
}
\]

for frequencies where \(\cos(\omega/2) > 0\).

%%%%%%%%%%%%%%%%%%%%%%%%%%%%%%%%%%%%%%%%%%%%%%%%%%%%%%%%%%%%%%%%%%%%%%%%%%%%%%%%

\section{Phase Delay}

Phase delay is defined as

\[
\boxed{
\tau_p(\omega)=
-\frac{\phi(\omega)}{\omega}
}
\]

It represents the delay experienced by a single sinusoidal component.

%%%%%%%%%%%%%%%%%%%%%%%%%%%%%%%%%%%%%%%%%%%%%%%%%%%%%%%%%%%%%%%%%%%%%%%%%%%%%%%%

\subsection{Example}

For

\[
\phi(\omega)=-\frac{\omega}{2},
\]

\[
\tau_p(\omega)=
-\frac{-\omega/2}{\omega}
=
\frac12.
\]

\[
\boxed{
\tau_p(\omega)=0.5\text{ samples}
}
\]

%%%%%%%%%%%%%%%%%%%%%%%%%%%%%%%%%%%%%%%%%%%%%%%%%%%%%%%%%%%%%%%%%%%%%%%%%%%%%%%%

\section{Group Delay}

Group delay is

\[
\boxed{
\tau_g(\omega)=
-\frac{d\phi(\omega)}{d\omega}
}
\]

It represents the delay experienced by a narrowband group of frequencies (an
envelope).

%%%%%%%%%%%%%%%%%%%%%%%%%%%%%%%%%%%%%%%%%%%%%%%%%%%%%%%%%%%%%%%%%%%%%%%%%%%%%%%%

\subsection{Example}

Again,

\[
\phi(\omega)=-\frac{\omega}{2}.
\]

Then

\[
\tau_g(\omega)=
-\frac{d}{d\omega}\left(-\frac{\omega}{2}\right)
=
\frac12.
\]

\[
\boxed{
\tau_g(\omega)=0.5\text{ samples}
}
\]

%%%%%%%%%%%%%%%%%%%%%%%%%%%%%%%%%%%%%%%%%%%%%%%%%%%%%%%%%%%%%%%%%%%%%%%%%%%%%%%%

\section{Phase Delay vs Group Delay}

\begin{table}[H]
\centering
\caption{Comparison of phase and group delay}
\begin{tabular}{|p{5cm}|p{4cm}|p{4cm}|}
\hline
\textbf{Quantity} & \textbf{Formula} & \textbf{Meaning}\\
\hline
Phase delay & \(-\phi/\omega\) & Delay of a sinusoid\\
\hline
Group delay & \(-d\phi/d\omega\) & Delay of an envelope\\
\hline
\end{tabular}
\end{table}

For a linear-phase system, both are constant and equal.

%%%%%%%%%%%%%%%%%%%%%%%%%%%%%%%%%%%%%%%%%%%%%%%%%%%%%%%%%%%%%%%%%%%%%%%%%%%%%%%%

\section{Why Group Delay Matters}

A digital communication signal contains many frequency components.

If different frequencies experience different delays,

\begin{itemize}
\item the waveform shape changes,
\item pulses spread,
\item intersymbol interference occurs.
\end{itemize}

Therefore, communication systems prefer filters with \textbf{constant group delay}.

%%%%%%%%%%%%%%%%%%%%%%%%%%%%%%%%%%%%%%%%%%%%%%%%%%%%%%%%%%%%%%%%%%%%%%%%%%%%%%%%

\section{Linear-Phase Systems}

A system has linear phase if

\[
\boxed{
\phi(\omega)=-(\alpha\omega+\beta)
}
\]

where \(\alpha\) and \(\beta\) are constants.

Then

\[
\boxed{
\tau_g(\omega)=\alpha
}
\]

which is independent of frequency.

%%%%%%%%%%%%%%%%%%%%%%%%%%%%%%%%%%%%%%%%%%%%%%%%%%%%%%%%%%%%%%%%%%%%%%%%%%%%%%%%

\section{Symmetric FIR Condition}

A real FIR filter has linear phase if

\[
\boxed{
h[n]=h[N-1-n]
}
\]

or

\[
\boxed{
h[n]=-h[N-1-n].
}
\]

%%%%%%%%%%%%%%%%%%%%%%%%%%%%%%%%%%%%%%%%%%%%%%%%%%%%%%%%%%%%%%%%%%%%%%%%%%%%%%%%

\section{Example: Symmetric FIR}

\[
h[n]=\{1,2,1\}.
\]

Check symmetry:

\[
h[0]=h[2]=1.
\]

Hence the filter is symmetric and therefore linear phase.

%%%%%%%%%%%%%%%%%%%%%%%%%%%%%%%%%%%%%%%%%%%%%%%%%%%%%%%%%%%%%%%%%%%%%%%%%%%%%%%%

\section{Frequency Response}

\[
H(e^{j\omega})=
1+2e^{-j\omega}+e^{-j2\omega}.
\]

Factorize:

\[
H(e^{j\omega})=
e^{-j\omega}(2+2\cos\omega).
\]

Thus,

\[
\boxed{
\phi(\omega)=-\omega
}
\]

and

\[
\boxed{
\tau_g(\omega)=1
}
\]

sample.

%%%%%%%%%%%%%%%%%%%%%%%%%%%%%%%%%%%%%%%%%%%%%%%%%%%%%%%%%%%%%%%%%%%%%%%%%%%%%%%%

\section{The Four Linear-Phase FIR Types}

%%%%%%%%%%%%%%%%%%%%%%%%%%%%%%%%%%%%%%%%%%%%%%%%%%%%%%%%%%%%%%%%%%%%%%%%%%%%%%%%

\subsection{Type I}

\begin{itemize}
\item Symmetric
\item Odd length
\item Most commonly used
\end{itemize}

%%%%%%%%%%%%%%%%%%%%%%%%%%%%%%%%%%%%%%%%%%%%%%%%%%%%%%%%%%%%%%%%%%%%%%%%%%%%%%%%

\subsection{Type II}

\begin{itemize}
\item Symmetric
\item Even length
\item Zero at \(\omega=\pi\)
\end{itemize}

%%%%%%%%%%%%%%%%%%%%%%%%%%%%%%%%%%%%%%%%%%%%%%%%%%%%%%%%%%%%%%%%%%%%%%%%%%%%%%%%

\subsection{Type III}

\begin{itemize}
\item Antisymmetric
\item Odd length
\item Zero at \(\omega=0\)
\end{itemize}

%%%%%%%%%%%%%%%%%%%%%%%%%%%%%%%%%%%%%%%%%%%%%%%%%%%%%%%%%%%%%%%%%%%%%%%%%%%%%%%%

\subsection{Type IV}

\begin{itemize}
\item Antisymmetric
\item Even length
\item Zeros at both \(\omega=0\) and \(\omega=\pi\)
\end{itemize}

%%%%%%%%%%%%%%%%%%%%%%%%%%%%%%%%%%%%%%%%%%%%%%%%%%%%%%%%%%%%%%%%%%%%%%%%%%%%%%%%

\section{Type I Example}

\[
h[n]=\{1,3,3,1,0\}.
\]

Length:

\[
N=5 \quad (\text{odd})
\]

and the sequence is symmetric.

Hence,

\[
\boxed{
\text{Type I linear-phase FIR}
}
\]

%%%%%%%%%%%%%%%%%%%%%%%%%%%%%%%%%%%%%%%%%%%%%%%%%%%%%%%%%%%%%%%%%%%%%%%%%%%%%%%%

\section{Group Delay of Linear-Phase FIR}

For an \(N\)-tap linear-phase FIR filter,

\[
\boxed{
\tau_g=
\frac{N-1}{2}
}
\]

%%%%%%%%%%%%%%%%%%%%%%%%%%%%%%%%%%%%%%%%%%%%%%%%%%%%%%%%%%%%%%%%%%%%%%%%%%%%%%%%

\subsection{Examples}

\[
N=5 \Rightarrow \tau_g=2
\]

\[
N=8 \Rightarrow \tau_g=3.5
\]

%%%%%%%%%%%%%%%%%%%%%%%%%%%%%%%%%%%%%%%%%%%%%%%%%%%%%%%%%%%%%%%%%%%%%%%%%%%%%%%%

\section{Phase Unwrapping}

The phase returned by \(\tan^{-1}\) is restricted to

\[
(-\pi,\pi].
\]

When the true phase crosses \(\pm\pi\), discontinuities appear.

%%%%%%%%%%%%%%%%%%%%%%%%%%%%%%%%%%%%%%%%%%%%%%%%%%%%%%%%%%%%%%%%%%%%%%%%%%%%%%%%

\subsection{Wrapped Phase}

\[
170^\circ,\;175^\circ,\;-179^\circ,\;-175^\circ
\]

%%%%%%%%%%%%%%%%%%%%%%%%%%%%%%%%%%%%%%%%%%%%%%%%%%%%%%%%%%%%%%%%%%%%%%%%%%%%%%%%

\subsection{Unwrapped Phase}

\[
170^\circ,\;175^\circ,\;181^\circ,\;185^\circ
\]

Unwrapping adds multiples of \(2\pi\) to maintain continuity.

%%%%%%%%%%%%%%%%%%%%%%%%%%%%%%%%%%%%%%%%%%%%%%%%%%%%%%%%%%%%%%%%%%%%%%%%%%%%%%%%

\section{Distortionless Transmission}

A system is distortionless if

\[
\boxed{
y[n]=Kx[n-n_0]
}
\]

for some constant gain \(K\) and delay \(n_0\).

In the frequency domain,

\[
\boxed{
H(e^{j\omega})=
Ke^{-j\omega n_0}
}
\]

Hence,

\begin{itemize}
\item magnitude is constant,
\item phase is linear,
\item group delay is constant.
\end{itemize}

%%%%%%%%%%%%%%%%%%%%%%%%%%%%%%%%%%%%%%%%%%%%%%%%%%%%%%%%%%%%%%%%%%%%%%%%%%%%%%%%

\section{All-Pass Systems}

An all-pass filter has

\[
\boxed{
|H(e^{j\omega})|=1
}
\]

for all frequencies.

It changes only the phase.

%%%%%%%%%%%%%%%%%%%%%%%%%%%%%%%%%%%%%%%%%%%%%%%%%%%%%%%%%%%%%%%%%%%%%%%%%%%%%%%%

\section{Example}

\[
H(z)=
\frac{a+z^{-1}}{1+az^{-1}},
\qquad |a|<1.
\]

This is a first-order all-pass filter.

%%%%%%%%%%%%%%%%%%%%%%%%%%%%%%%%%%%%%%%%%%%%%%%%%%%%%%%%%%%%%%%%%%%%%%%%%%%%%%%%

\section{Magnitude Verification}

Substituting \(z=e^{j\omega}\),

\[
|H(e^{j\omega})|=1.
\]

The filter is useful for phase equalization and delay compensation.

%%%%%%%%%%%%%%%%%%%%%%%%%%%%%%%%%%%%%%%%%%%%%%%%%%%%%%%%%%%%%%%%%%%%%%%%%%%%%%%%

\section{Minimum-Phase Systems}

A system is minimum phase if

\begin{itemize}
\item all poles are inside the unit circle,
\item all zeros are also inside the unit circle.
\end{itemize}

Such systems have the smallest possible phase delay for a given magnitude response.

%%%%%%%%%%%%%%%%%%%%%%%%%%%%%%%%%%%%%%%%%%%%%%%%%%%%%%%%%%%%%%%%%%%%%%%%%%%%%%%%

\section{Maximum-Phase Systems}

If one or more zeros lie outside the unit circle, the system is maximum phase.

These systems have larger phase delay.

%%%%%%%%%%%%%%%%%%%%%%%%%%%%%%%%%%%%%%%%%%%%%%%%%%%%%%%%%%%%%%%%%%%%%%%%%%%%%%%%

\section{Mixed-Phase Systems}

If some zeros are inside and some are outside the unit circle, the system is mixed
phase.

Most arbitrary FIR filters are mixed phase unless designed specifically for linear
or minimum phase.

%%%%%%%%%%%%%%%%%%%%%%%%%%%%%%%%%%%%%%%%%%%%%%%%%%%%%%%%%%%%%%%%%%%%%%%%%%%%%%%%

\section{Solved Example 12.5}

Find the phase and group delay of

\[
H(z)=1+z^{-1}.
\]

%%%%%%%%%%%%%%%%%%%%%%%%%%%%%%%%%%%%%%%%%%%%%%%%%%%%%%%%%%%%%%%%%%%%%%%%%%%%%%%%

\subsection{Solution}

\[
H(e^{j\omega})=
2e^{-j\omega/2}\cos\frac{\omega}{2}.
\]

Therefore,

\[
\phi(\omega)=-\frac{\omega}{2}.
\]

Group delay:

\[
\tau_g=
-\frac{d}{d\omega}\left(-\frac{\omega}{2}\right)
=
\frac12.
\]

\[
\boxed{
\tau_g=0.5\text{ samples}
}
\]

%%%%%%%%%%%%%%%%%%%%%%%%%%%%%%%%%%%%%%%%%%%%%%%%%%%%%%%%%%%%%%%%%%%%%%%%%%%%%%%%

\section{Solved Example 12.6}

Determine whether

\[
h[n]=\{1,2,3,2,1\}
\]

has linear phase.

%%%%%%%%%%%%%%%%%%%%%%%%%%%%%%%%%%%%%%%%%%%%%%%%%%%%%%%%%%%%%%%%%%%%%%%%%%%%%%%%

\subsection{Check Symmetry}

\[
h[0]=h[4]=1,
\qquad
h[1]=h[3]=2.
\]

The sequence is symmetric.

\[
\boxed{
\text{Linear-phase FIR}
}
\]

Length:

\[
N=5.
\]

Hence,

\[
\boxed{
\tau_g=\frac{5-1}{2}=2
}
\]

samples.

%%%%%%%%%%%%%%%%%%%%%%%%%%%%%%%%%%%%%%%%%%%%%%%%%%%%%%%%%%%%%%%%%%%%%%%%%%%%%%%%

\section{Solved Example 12.7}

Find the group delay of an 8-tap linear-phase FIR filter.

%%%%%%%%%%%%%%%%%%%%%%%%%%%%%%%%%%%%%%%%%%%%%%%%%%%%%%%%%%%%%%%%%%%%%%%%%%%%%%%%

\subsection{Solution}

\[
\tau_g=
\frac{8-1}{2}
=
3.5.
\]

\[
\boxed{
3.5\text{ samples}
}
\]

%%%%%%%%%%%%%%%%%%%%%%%%%%%%%%%%%%%%%%%%%%%%%%%%%%%%%%%%%%%%%%%%%%%%%%%%%%%%%%%%

\section{Solved Example 12.8}

A system has

\[
H(e^{j\omega})=
3e^{-j2\omega}.
\]

Determine whether it is distortionless.

%%%%%%%%%%%%%%%%%%%%%%%%%%%%%%%%%%%%%%%%%%%%%%%%%%%%%%%%%%%%%%%%%%%%%%%%%%%%%%%%

\subsection{Solution}

Magnitude:

\[
|H|=3 \quad (\text{constant})
\]

Phase:

\[
\phi(\omega)=-2\omega
\]

which is linear.

Therefore,

\[
\boxed{
\text{The system is distortionless}
}
\]

with gain \(3\) and delay \(2\) samples.

%%%%%%%%%%%%%%%%%%%%%%%%%%%%%%%%%%%%%%%%%%%%%%%%%%%%%%%%%%%%%%%%%%%%%%%%%%%%%%%%

\section{Solved Example 12.9}

Identify the FIR type:

\[
h[n]=\{1,0,-1\}.
\]

Check antisymmetry:

\[
h[0]=-h[2].
\]

Length is odd (\(N=3\)).

\[
\boxed{
\text{Type III linear-phase FIR}
}
\]

%%%%%%%%%%%%%%%%%%%%%%%%%%%%%%%%%%%%%%%%%%%%%%%%%%%%%%%%%%%%%%%%%%%%%%%%%%%%%%%%

\section{Important GATE Points}

\begin{enumerate}
\item Frequency response is evaluated on the unit circle.
\item Group delay is the derivative of phase.
\item Constant group delay implies no waveform distortion.
\item Symmetric or antisymmetric FIR filters have linear phase.
\item Group delay of an \(N\)-tap linear-phase FIR is \((N-1)/2\).
\item All-pass filters have unity magnitude.
\item Minimum-phase systems have all poles and zeros inside the unit circle.
\end{enumerate}

%%%%%%%%%%%%%%%%%%%%%%%%%%%%%%%%%%%%%%%%%%%%%%%%%%%%%%%%%%%%%%%%%%%%%%%%%%%%%%%%

\section{Chapter Summary}

In this chapter, we studied the phase response of discrete-time systems and
defined phase delay and group delay.

We established the conditions for linear-phase FIR filters, classified the four
FIR linear-phase types, and discussed phase unwrapping, all-pass filters,
minimum-phase systems and distortionless transmission.

These concepts are fundamental in digital communication, audio processing and
high-performance filter design.

%%%%%%%%%%%%%%%%%%%%%%%%%%%%%%%%%%%%%%%%%%%%%%%%%%%%%%%%%%%%%%%%%%%%%%%%%%%%%%%%

\begin{tcolorbox}[title=One-Minute Revision,colback=blue!5]

\[
H(e^{j\omega})=
|H(e^{j\omega})|e^{j\phi(\omega)}
\]

\[
\tau_p(\omega)=
-\frac{\phi(\omega)}{\omega}
\]

\[
\tau_g(\omega)=
-\frac{d\phi(\omega)}{d\omega}
\]

Linear phase:

\[
\phi(\omega)=-(\alpha\omega+\beta)
\]

Symmetric FIR:

\[
h[n]=h[N-1-n]
\]

Antisymmetric FIR:

\[
h[n]=-h[N-1-n]
\]

Group delay:

\[
\tau_g=\frac{N-1}{2}
\]

Distortionless:

\[
H(e^{j\omega})=Ke^{-j\omega n_0}
\]

All-pass:

\[
|H(e^{j\omega})|=1
\]

Minimum phase:

\[
|p_i|<1,\quad |z_i|<1
\]

\end{tcolorbox}
%%%%%%%%%%%%%%%%%%%%%%%%%%%%%%%%%%%%%%%%%%%%%%%%%%%%%%%%%%%%%%%%%%%%%%%%%%%%%%%%